\chapter{Installing and starting}\label{ch:installing}

\section{What you need}
\begin{itemize}
  \item A Windows PC (or Linux, see below). A screen of 1024 by 768 pixels is enough.
  \item For the 3D view: a graphics card that can do OpenGL 2. Almost every PC from 2009 on can. Without it, the 3D button is grey and the 2D picture does everything else.
  \item To run the G-code: LinuxCNC on the lathe, and a way to carry the file there.
\end{itemize}

\section{Installing}
WOP for Windows comes as a zip file. Unpack it anywhere on your PC, for example to \texttt{C:\textbackslash WOP}. Nothing is installed. Start only the \file{WOP.exe} that is in the unpacked folder; do not move the file out of it.

The program is not signed. The first time, Windows (SmartScreen) or a virus scanner may warn you. If you got the zip from the place named in the release notes, you can allow it.

\begin{notebox}
On Linux, WOP is installed with a script, and the rest of this manual applies as it is. Developers can also run WOP from the source code with Python. Both are described in the file \file{README.md}.
\end{notebox}

\section{Starting WOP}
Double click \file{WOP.exe}. A window opens with an empty project. Use \menupath{File}{Open...} to open a project file. The example projects in the folder \file{examples} show how finished projects look. They carry the machine of their author (the tool change position and the limits of his lathe): choose your own machine before you export G-code from one of them.

The first time, look at \menupath{File}{Machines...} and make a machine for your lathe (Chapter~\ref{ch:machines}). Then choose it in the Setup.

\section{Where WOP keeps its files}
\subsection{The user folder}
Your settings, machines, libraries and postprocessors are not in the project. They are in the \emph{user folder}. WOP uses one of two places:
\begin{itemize}
  \item If there is a folder \file{user_data} next to \file{WOP.exe}, WOP uses it. Then you can copy the whole program folder to another PC and take everything with you.
  \item Otherwise it uses \texttt{\%APPDATA\%\textbackslash WOP}, in the folder of your Windows user.
\end{itemize}
To change from one to the other, make or remove the folder \file{user_data}. Files are not moved by themselves.

\subsection{Libraries and settings}
The user folder holds the settings (\file{settings.json}), your machines, the libraries of tools, workholding and materials, your own postprocessors, and two log files. Chapter~\ref{ch:app_files} lists them.
